A integração com o barramento CAN foi encapsulada em uma camada baseada no driver
TWAI da Espressif, que fornece a implementação de CAN clássico no ESP-IDF
\cite{espressif_twai_programming_guide}. Essa parte foi validada em bancada com
transceptores VP230, mas não foi o foco da campanha principal. Como a montagem em
protoboard introduzia falhas de contato, terminação e alimentação, os testes finais
isolaram a camada WCAN, que é a contribuição central deste trabalho.

\section{Arquitetura do WCAN}

O WCAN foi estruturado como um protocolo de aplicação sobre ESP-NOW. A interface
\texttt{wcan\_init} abstrai o transporte sem fio e permite que a aplicação trate os
dados recebidos como mensagens com semântica CAN. A camada empacota amostras por
identificador CAN, aplica filtros no receptor, envia os pacotes pelo rádio e entrega
os dados aceitos ao callback da aplicação.

Para suprir a necessidade de se investigar a viabilidade entre as abordagens citadas
anteriormente, foram implementadas duas variantes de transporte para o WCAN, broadcast 
e multicast, selecionadas em tempo de compilação. As duas usam o mesmo formato de pacote 
e a mesma aplicação de teste, o que permitiu comparar o impacto dessas estratégias de
endereçamento de forma isolada na confiabilidade sem alterar a instrumentação. O papel de 
cada placa, sensor, receptor ou nó ocioso, é definido no boot por uma linha enviada pela 
serial, junto com frequência de geração, CAN ID base, quantidade de IDs e valor de \texttt{linger\_ms}.

\section{Formato dos Pacotes}

A unidade de comunicação é a estrutura \texttt{data\_packet\_t}, composta pelo MAC
de origem, CAN ID, contador de ticks, quantidade de amostras e vetor de dados de 32
bits. Para transmissão, o pacote é serializado com 9 bytes de cabeçalho: 4 bytes de
CAN ID, 4 bytes de contador e 1 byte de quantidade de amostras. Como o ESP-NOW
clássico permite até 250 bytes de payload, restam 241 bytes para dados, o que limita
cada lote a 60 amostras de 4 bytes \cite{espressif_espnow_programming_guide}.

Os identificadores de controle foram colocados fora do espaço válido de IDs CAN
estendidos. Como o CAN estendido usa 29 bits, com valor máximo \texttt{0x1FFFFFFF},
a faixa \texttt{0xE0000000} foi reservada para sentinelas internas do WCAN
\cite{iso11898_1_2024}. No broadcast, esse valor identifica ACKs. No multicast,
identifica mensagens de descoberta, chamadas \texttt{HELLO}. Como apenas uma
variante é compilada por vez, não há conflito no firmware final.

\section{Filas e Agrupamento de Amostras}

Os callbacks do ESP-NOW foram mantidos curtos, conforme a recomendação da Espressif
de deslocar operações demoradas para filas e tasks de menor prioridade
\cite{espressif_espnow_programming_guide}. Na recepção, o callback apenas copia e
valida o pacote recebido. Filtragem, deduplicação, chamada da aplicação,
descoberta e retransmissões são tratadas fora desse contexto.

No caminho de transmissão, cada CAN ID possui uma fila própria. A task associada
consome amostras, espera até \texttt{linger\_ms} e fecha o lote quando a janela
expira ou quando 60 amostras são acumuladas. Esse batching é essencial porque o
ESP-NOW é um enlace de pacotes, com custo fixo por quadro. Para aquisição de dados
com frequência conhecida, o atraso controlado por \texttt{linger\_ms} reduz o uso do
canal e a pressão sobre filas. Quando \texttt{linger\_ms = 0}, o sistema envia uma
amostra por pacote, configuração mantida apenas como experimento de limite. Estudos
e medições práticas de ESP-NOW também indicam que tamanho de pacote, topologia e
condições do enlace afetam diretamente vazão e atraso \cite{hackaday_espnow_throughput,urazayev2023_indoor_espnow,wedyanti2025_espnow_multinode}.

\section{Transportes Broadcast e Multicast}

No broadcast, todos os pacotes de dados são enviados para
\texttt{FF:FF:FF:FF:FF:FF}. Os receptores no mesmo canal aplicam seus filtros por
CAN ID e, se aceitarem o pacote, enviam um ACK direcionado ao sensor. O transmissor
usa o par CAN ID e contador para correlacionar confirmações, retransmitindo quando
nenhum ACK válido chega dentro da janela esperada. Quando existem múltiplos
receptores, o primeiro ACK libera o próximo lote. Essa escolha replica a semântica
do ACK dominante do CAN, na qual o transmissor sabe que algum nó recebeu a
mensagem, mas não identifica quais nós a receberam \cite{iso11898_1_2024}.

No multicast, o sensor mantém uma tabela de assinantes interessados em cada CAN ID.
Enquanto não há assinantes vivos, o sensor transmite em broadcast para permitir
descoberta imediata. Um receptor interessado responde com \texttt{HELLO}, e o sensor
passa a incluí-lo nas próximas transmissões multicast. Receptores também enviam
\texttt{HELLO} periodicamente, permitindo que novos nós sejam descobertos mesmo após
o início da transmissão. O envio multicast usa o retorno do driver ESP-NOW como
sinal de sucesso no nível MAC, mas esse retorno não substitui uma confirmação
explícita de processamento pela aplicação \cite{espressif_espnow_programming_guide}.

As duas estratégias refletem compromissos diferentes. O broadcast é simples e
funciona sem descoberta prévia, mas só garante que ao menos um receptor interessado
confirmou o pacote. O multicast é o caminho natural para múltiplos receptores, mas
depende de descoberta, liveness e gerenciamento de memória mais robustos. Por isso,
a decisão entre as variantes foi tratada como questão experimental.

\section{Integração com o CAN Existente}

A integração preserva o CAN ID original. O receptor WCAN entrega um
\texttt{data\_packet\_t} à aplicação, que pode retransmitir os dados no barramento
físico via TWAI usando o mesmo identificador. Assim, os demais módulos do veículo
podem tratar os dados recebidos por rádio como mensagens CAN comuns
\cite{silva2009_rede_can_formula_sae,blatt2020_aquisicao_formula_sae}. Nos testes
quantitativos, o callback registrou os pacotes em log serial para análise de entrega
e latência. No uso embarcado, essa fronteira é o ponto de injeção no barramento CAN
local.
